\clearpage
\begin{summarybox}
\textbf{C087｜双曲线焦点弦长公式与取值范围——快记／快懂／快用}

\textbf{快记（两个公式，两类取值范围）：}
\[
\boxed{L=\frac{2ab^2(1+k^2)}{|a^2k^2-b^2|}\quad(k\ne\pm b/a)}
\]
\[
\boxed{L=\frac{2ab^2}{|a^2-c^2\cos^2\theta|}
\quad\left(\cos\theta\ne\pm\frac ac\right)}
\]
\[
\boxed{\begin{array}{rl}
\text{异支}~(|k|<b/a):&L\in[2a,+\infty),\ \text{水平取等};\\[2pt]
\text{同支}~(|k|>b/a\text{或竖直}):&
L\in[2b^2/a,+\infty),\ \text{竖直取等}.
\end{array}}
\]

\textbf{快懂（最小认知负担）：}
\textbf{跨两支}\ $\Rightarrow$ 先跨越实轴两顶点，最短 $2a$；
\textbf{同一支}\ $\Rightarrow$ 过焦点竖切，最短 $2b^2/a$。
旋转到\textbf{渐近线方向}附近，远端交点跑向无穷，弦长越来越大。

\textbf{快用（考场选择）：}
\begin{itemize}
\item \textbf{最快证明：}设 $t=x-c$，二次方程 $At^2+Bt+C=0$，
$\Delta=4a^2b^4(1+k^2)$，由根差求弦长；
\item \textbf{最低门槛：}异支用 $L\ge\Delta x\ge2a$；
同支把 $L$ 写成 $2b^2/a$ 加一个正项，补上竖直取等；
\item \textbf{常规解析：}联立 $y=k(x-c)$ 与双曲线，根与系数关系＋弦长公式；
\item \textbf{几何补充：}右支点的焦半径 $r=b^2/(a-c\cos\varphi)$，
正反向求和，再用 $1/(a-v)+1/(a+v)\ge2/a$。
\end{itemize}

\textbf{三重防错：}斜率 $\pm b/a$ 不能用；\textbf{竖直线必须另算}；
判断同支/异支看 $|k|$，不要看 $k$ 的符号。

\textbf{与相邻结论区分：}本结论只研究\textbf{过焦点的弦长}，
不把教材后面的“渐近线面积、交点性质”混进 C087。
\end{summarybox}
